% AI-written, AI cross-reviewed; no human sign-off; not a claim that RH is proved
% Atlas unit P034. Local manuscript v3, item 7.11, heading on printed PDF p. 79.
% PositivityTheory.KB.bOdd_zk
% paper2/inputs_one.tex:476–533. Verbatim excerpt; SHA-256 09a233150c7acb8f8e798d9658f9cf31965028ba23d9c66d6007467b358056c0.
% Original TeX uses the macros and cross-reference keys of the source manuscript.


% Start of atlas edition annotation. The following block is outside the original excerpt.
\paragraph{Status of named inputs in this reading.}
\noindent\nolinkurl{PositivityTheory.KB.bStarOdd} --- in-repo constant (\path{Development/AutomorphicArithmetic/AutomorphicVarieties/LiteratureResults/KTheoryStability/ReductionBoundsTorsion/FinalOdd.lean:8--11}); not compared.\par
% End of atlas edition annotation.

\begin{formalstatement}[Odd homology after successive block operations]\label{pn:P034}
Let \(L\) be totally complex, \(k\ge0\) even, \(q=k+1<r_2(L)\), and
\(n\ge2q+4\). Fix any integer \(t>0\). Starting with \(z_0(x)=x\)
in \(H_q(G_n,\mathbb Z)\), define
\[
 z_{j+1}(x)=(\operatorname{diag}_t)_*z_j(x)
                 -t^{j+2}\operatorname{stab}_*z_j(x),
\]
where block repetition and stabilization both land in rank \(t^{j+1}n\).
The integral span of all \(z_k(x)\) has rank at most \(r_2(L)\).
\end{formalstatement}\nodestatus{P034}
\begin{outline}
Use the comparison in Blueprint Z2.K4~\bploc{r7:review:Z2.K4}. Choose a rank \(N\) above
\(t^kn\) at which Borel's stable real comparison represents classes
through degree \(q\) on \(\operatorname{Res}_{L/\mathbb Q}\mathrm{SL}_N\)
by invariant forms. The central directions are supplied by
\(d\log\det\) at infinity: scalar units modulo torsion span the
sum-zero real hyperplane by Dirichlet's theorem. The determinant-one
subgroup together with a torsion-free subgroup of scalar units has
finite index. Künneth gives the mixed products on this subgroup,
and transfer is injective on restriction with real coefficients.
Integration over straight simplices in the positive Hermitian cone
computes their pairing with real group homology.

An invariant alternating tensor is a combination of products of
alternating trace words, with a place attached to each trace.
One sees this by complexifying at the identity and contracting tensor
indices along the cycles of permutations. A product with \(j\)
trace factors scales by \(t^j\) under block repetition, and
stabilization preserves its evaluation. When \(2\le j\le k+1\) its
evaluation on \(z_k\) therefore vanishes: the step indexed \(j-2\)
subtracts exactly the same multiple. In odd degree the one-trace
evaluations have at most one real generator at each place. Indeed
reversing the word conjugates its trace on Hermitian matrices; after
alternation the real or imaginary part vanishes according to the
sign of reversal. All one-cycle permutations give the same remaining
trace up to sign.

Stabilization from \(t^kn\) to \(N\) is an integral homology isomorphism
in degree \(q\). The \(r_2(L)\) remaining evaluations separate the real
span of \(z_k(x)\): a vector in their common kernel pairs trivially
with every class by the preceding decomposition. Finally,
extension to \(\mathbb R\) commutes with the bar differential, block
repetition and stabilization. The rank of the original image is its
dimension after this extension.
\end{outline}

Borel's comparison is used at an auxiliary rank. His Proposition~12.2
computes the real ranks of the stable \(K_i\) of an order for \(i\ge2\);
for a totally complex field these ranks are \(r_2(L)\) in odd degrees
\(i\ge3\). The project's restriction \(k+1<r_2(L)\) belongs to its
stable-realization assertion. Borel's rank formula has no such
restriction and does not cover \(k=0\) (\(i=1\)); the central-unit
calculation above handles that degree. Transfer to the specified rank
and projection by the block recurrence are additional comparisons.
See A.~Borel, \emph{Stable real cohomology of arithmetic groups},
Theorems~7.5 and 11.1, \S\S9.5, 11.2--11.4, and Proposition~12.2.
\cite{r21-borel}
